\subsection{Tier B primary matrix}\label{sec:tierb}
The primary matrix comprises \val{JNRunsMain} mission runs, that is, \val{JNMain} held-out missions in
\val{JNConfigsMain} configurations with eight methods. Table~\ref{tab:main} lists utility, lost tasks and computation per
method, and Fig.~\ref{fig:main} relates utility to computation and shows the difference in every configuration.

Search pays off clearly under uncertain execution. Regret-2 repair reached a utility of \val{JURepair}\%, and
AnytimeSwarm-R reached \val{JURisk}\%, which is \val{JDmRepair} percentage points (pp) or \val{JURelRepair}\% more. The
full budget of 250\,ms was the weakest search method, with \val{JUFixed}\%, and AnytimeSwarm-R served \val{JDFixed}\,pp
more reward than it. The loss of the long budget stems from late plans. A task counts as lost to delay when it is part of
the computed plan but can no longer be served on time when the plan is released. With $\kappa_B=50$, the full budget
lost \val{JLostKFiftyFixed} tasks per mission in this way, against \val{JLostKFiftyRisk} for AnytimeSwarm-R, that is,
\val{JLostRedFixed}\% fewer. Without its delay cost, AnytimeSwarm lost \val{JLostKFiftyNoDelay} tasks per mission in this
way, but its utility was not lower (\val{JUNoDelay}\% against \val{JUAS}\%). The delay cost thus prevents losses from
late plans without a measurable gain in utility on this matrix.

Among the remaining search methods, the differences are small. All seven search methods lie within
\val{JUSpreadSearch}\,pp of each other. AnytimeSwarm-R reached the highest mean utility with the least computation of the adaptive rules,
\val{JMsRisk}\,ms per call against \val{JMsAS}\,ms for AnytimeSwarm and \val{JMsStag}\,ms for the stagnation rule. The
reduction of \val{JMsRedAS}\% relative to AnytimeSwarm is significant (\val{JDMsAS}\,ms per call), so the second
hypothesis holds. The primary hypothesis does not hold at the level fixed in the protocol. The pooled advantage of
AnytimeSwarm-R was \val{JDFixedForty}\,pp over ALNS with 40\,ms, \val{JDAS}\,pp over AnytimeSwarm and \val{JDStag}\,pp
over the stagnation rule, with Holm-adjusted $p$-values of \val{JPFixedForty}, \val{JPAS} and \val{JPStag}. Only the
first interval excludes zero before correction. Per configuration and without correction for multiple testing,
AnytimeSwarm-R was better than ALNS with 40\,ms in \val{JWinFixedForty} of the \val{JNConfigsMain} configurations, better
than the stagnation rule in \val{JWinStag} and never different from AnytimeSwarm, which is consistent with chance
(Fig.~\ref{fig:main}b). The ablation that keeps only the latency term reached \val{JURiskZero}\%, close to AnytimeSwarm.
Together with the misspecification runs below, this suggests that a non-zero error spread mainly makes the rule stop
earlier, rather than that the calibration of the spread matters.

The fixed budget of 40\,ms is a strong reference, because it was matched to the mean computation of AnytimeSwarm on
development data. Its weakness lies in the allocation across fleet scales. AnytimeSwarm-R used \val{JMsScSmallRisk}\,ms
per call on the small scale and \val{JMsScLargeRisk}\,ms on the large scale. On the large scale, it served
\val{JDScLargeFixedForty}\,pp more reward than the budget of 40\,ms and was on a par with the full budget
(\val{JDScLargeFixed}\,pp). On the small scale, it served \val{JDScSmallFixed}\,pp more than the full budget and was on a
par with the budget of 40\,ms (\val{JDScSmallFixedForty}\,pp). A single fixed budget is therefore either too long for
small fleets or too short for large ones, whereas the adaptive rules find a suitable amount of computation on each scale
without tuning. The advantage over ALNS with 40\,ms was \val{JDFFKappa10}\,pp for $\kappa_B=10$ and
\val{JDFFKappa50}\,pp for $\kappa_B=50$, and \val{JDFFExecLow}\,pp under low and \val{JDFFExecMid}\,pp under mid
disturbances.

\inputtable{tab_main}

\begin{figure}[t]
\centering
\fig{fig_main}{\linewidth}
\caption{Tier B primary matrix, held-out seeds 3000--3029: (a) mean utility against mean planning time per call over
all \val{JNConfigsMain} configurations (regret-2 repair, at \val{JURepair}\% and 1\,ms, is not shown); (b) paired
difference of AnytimeSwarm-R minus ALNS with 40\,ms and minus AnytimeSwarm per configuration (execution level, latency
level, $\kappa_B$, trigger) with 95\% bootstrap intervals over \val{JNMain} missions.}\label{fig:main}
\end{figure}

\subsection{Replanning triggers}\label{sec:res-triggers}
The trigger mattered more than the stopping rule. Replanning every 60\,s in addition to the events raised the number of
planning calls from \val{JCallsEvent} to \val{JCallsPeriodic} per mission and the utility of every method, by
\val{JTrigMin}\,pp to \val{JTrigMax}\,pp, with all intervals above zero (\val{JTrigRisk}\,pp for AnytimeSwarm-R,
\val{JTrigFixedForty}\,pp for ALNS with 40\,ms). The third hypothesis therefore holds. Periodic replanning probably helps
because it corrects the accumulated deviations of the executed routes, which the event trigger only notices at the next
event. The advantage of AnytimeSwarm-R over ALNS with 40\,ms was \val{JDFixedFortyEvent}\,pp with the event trigger and
\val{JDFixedFortyPeriodic}\,pp with the periodic trigger. The two values are close, and the protocol did not test their
difference.

\subsection{Scalability to 40 UAVs}\label{sec:res-xl}
With 40 UAVs and 280 tasks, \val{JNXL} held-out missions were run under mid disturbances in \val{JXLNConf}
configurations of latency level and latency factor. All five methods reached between \val{JXLUAS}\% and
\val{JXLUStag}\%. AnytimeSwarm-R reached \val{JXLURisk}\% with \val{JXLMsRisk}\,ms per call. Its pooled difference was
\val{JXLDFixedOneEighty}\,pp to ALNS with 180\,ms, \val{JXLDAS}\,pp to AnytimeSwarm and \val{JXLDStag}\,pp to the
stagnation rule. After Holm correction over the four comparisons of the planned analysis, none of them is
significant ($p_{\mathrm{Holm}}=$ \val{JXLPAS} for AnytimeSwarm), so the fifth hypothesis is not supported either.
AnytimeSwarm-R again needed less computation than AnytimeSwarm (\val{JXLDMsAS}\,ms per call). The stagnation rule
stopped earliest, after \val{JXLMsStag}\,ms, and was significantly better than AnytimeSwarm-R in \val{JXLLossStag} of the
\val{JXLNConf} configurations, which had no latency and $\kappa_B=50$. On this scale, the rules used between \val{JXLMsStag} and 250\,ms
per call, and utility hardly changed within this range of computation. The equal-computation advantage that
AnytimeSwarm showed with 40 UAVs under exact execution (Table~\ref{tab:scale}) therefore did not translate into a
utility gain under uncertain execution.

\subsection{Sensitivity analyses}\label{sec:res-sens}
The sensitivity runs use the held-out missions of the primary matrix without latency and are descriptive. They vary the
execution level beyond the primary matrix, the latency factor and the trigger in \val{JNSensConf} configurations, and
the error model in \val{JNMisConf} further configurations.

Under high disturbances, which lie outside the range of the primary matrix but within that of the fitted error models,
utility fell to about \val{JSensUHighRisk}\% for all methods. The combination of high disturbances and a slow planner
($\kappa_B=50$) gave the largest advantage of the risk-aware cost in this study. AnytimeSwarm-R served
\val{JSensHighKFiftyAS}\,pp more reward than AnytimeSwarm and used \val{JSensMsHighKFiftyRisk}\,ms instead of
\val{JSensMsHighKFiftyAS}\,ms per call. Its difference to ALNS with 40\,ms was \val{JSensHighKFiftyFixedForty}\,pp and to
the stagnation rule \val{JSensHighKFiftyStag}\,pp, so the gain over AnytimeSwarm brought it to the level of the
baselines rather than beyond them. With $\kappa_B=1$, planning time hardly delays the plan, both delay costs are close to
zero because they scale with the shift $\kappa\tau$, and the two rules behave alike. AnytimeSwarm-R then differed from
AnytimeSwarm by \val{JSensKOneAS}\,pp and \val{JSensKOneMsDiff}\,ms per call.

Replanning when a UAV deviated by more than 30\,s from its planned finish time raised the number of calls from
\val{JSensCallsEvent} to \val{JSensCallsDev} per mission and the utility by \val{JSensTrigDevRisk}\,pp for
AnytimeSwarm-R and \val{JSensTrigDevFixedForty}\,pp for ALNS with 40\,ms. Like the periodic trigger, it helped both
stopping rules by a similar amount. Under the deviation trigger, the difference of AnytimeSwarm-R was
\val{JSensDevFixedForty}\,pp to ALNS with 40\,ms and \val{JSensDevAS}\,pp to AnytimeSwarm.

The utility of AnytimeSwarm-R was insensitive to the error model. With the error model of the low, mid or high level
under low and mid disturbances, its difference to AnytimeSwarm ranged between \val{JMisAllLow}\,pp and
\val{JMisAllMid}\,pp, the correctly specified model was not better than a wrong one, and the mean utilities of all
versions lay within \val{JMisSpread}\,pp. Setting the error model to zero gave
\val{JMisAllZero}\,pp. The four misspecification configurations also belong to the primary matrix, so the correctly
specified method was run twice on the same missions. The two runs differ only in the measured planning times, because
the stopping rule reads a wall clock. They gave identical utility in \val{JRepeatSameShare}\% of the missions, and the
mean utility of a method differed by up to \val{JRepeatDiffMax}\,pp between the runs. Differences of this size between
stopping rules should therefore not be read as systematic.

\subsection{Tier C: PX4 software-in-the-loop with Gazebo}\label{sec:res-tierc}
All \val{JNCRuns} held-out flights of tier C were completed, that is, \val{JNC} missions with four methods each. Two runs
did not reach their altitude within 90\,s and were repeated as specified. A further \val{JCNInvalid} runs had a real-time
factor between \val{JCRTFInvalidMin} and \val{JCRTFInvalidMax}, below the validity range of 1.6 to 2.2 stated in the
protocol. The protocol did not specify how to treat such runs, so they were repeated once with the same seed, like a
technical failure, and the original runs were kept. The repeated runs reached the same utility as the original ones, and
the comparisons below are identical with either set (\val{JCDFixedFortyOrig}\,pp and \val{JCDASOrig}\,pp with the
original runs). The real-time factor of the analysed runs was \val{JCRTF} on average and at least \val{JCRTFmin}. A
released plan reached all vehicles after \val{JCLat}\,s of simulation time on average, which is included in the expected
latency of AnytimeSwarm-R.

Table~\ref{tab:tierc} summarises the results. AnytimeSwarm-R reached a utility of \val{JCURisk}\%, against
\val{JCUFixedForty}\% for ALNS with 40\,ms, \val{JCUAS}\% for AnytimeSwarm and \val{JCUFixed}\% for the full budget.
The paired differences were \val{JCDFixedForty}\,pp to ALNS with 40\,ms and \val{JCDAS}\,pp to AnytimeSwarm, with
Holm-adjusted $p$-values of \val{JCPFixedForty} and \val{JCPAS}. The fourth hypothesis is therefore not supported. With
\val{JNC} missions, the intervals are much wider than in tier B and cannot resolve differences of the size seen there.
Two effects of tier B reappear in tier C. The full budget again had the lowest mean utility, although its interval
includes zero (\val{JCDFixed}\,pp to AnytimeSwarm-R), and it was the only method that lost a noticeable number of tasks
to late plans,
\val{JCDropFixed} per mission against at most \val{JCDropAS} for the others. AnytimeSwarm-R again computed least among
the adaptive rules, with \val{JCMsRisk}\,ms per call against \val{JCMsAS}\,ms for AnytimeSwarm, a reduction of
\val{JCMsRedAS}\% (\val{JCDMsAS}\,ms).

\inputtable{tab_tierc}

The same \val{JNC} missions were also run in tier B with the rescaled parameters of tier C. Tier B predicted the outcome
of a flight well. The utility of a run in tier C correlated with that of the same run in tier B with $r=$ \val{JBCr}
(\val{JBCn} runs), and the mean utility per mission with $r=$ \val{JBCrMission} (Fig.~\ref{fig:bc}b). Tier C served
slightly more reward on average, \val{JBCUC}\% against \val{JBCUB}\%, so the kinematic model of tier B was, if anything,
slightly pessimistic for these missions. The small
differences between the methods were not predicted as well. The method means of the two tiers had a Kendall rank
correlation of \val{JBCTau}, and the sign of the per-mission difference between AnytimeSwarm-R and ALNS with 40\,ms
agreed in \val{JBCSignAgreeFixedForty}\% of the missions, counting ties as agreement only when both differences are zero.
Part of the per-run correlation is by construction, because both tiers draw the wind, speed and service disturbances
from the same random numbers, and much of it reflects the different difficulty of the missions. The correlation was
first computed with an erroneous pairing of runs in the analysis script, which was corrected after the results had
been seen; the hypothesis tests were not affected. In tier B, the advantage of AnytimeSwarm-R on these missions
was \val{JBgzDFixedForty}\,pp over ALNS with 40\,ms and \val{JBgzDAS}\,pp over AnytimeSwarm (Fig.~\ref{fig:bc}a). Both
tiers thus agree that the methods differ by at most a few percentage points on this scale, and that the full budget is
the weakest search method, while the order of the three short-computation methods is within the noise of both tiers.

\begin{figure}[t]
\centering
\fig{fig_bc}{\linewidth}
\caption{Tier B against tier C on the same \val{JNC} held-out missions (three UAVs, mid disturbances, $\kappa_B=50$):
(a) paired difference of AnytimeSwarm-R minus each baseline with 95\% bootstrap intervals; (b) utility of every run in
tier C against the same run in tier B (dashed: equality).}\label{fig:bc}
\end{figure}
